\subsection*{\texorpdfstring{Positivity on $\mathbb{Q}$}{Positivity on Q}}

\begin{definitionbox}{Definition}
\begin{definition}[Positive Rational]
\label{def:positive-rational}
A rational number $[a,b]\in\mathbb{Q}$ is \emph{positive} if
\[
a\cdot b>0_{\mathbb{Z}}.
\]
\end{definition}
\end{definitionbox}

\begin{remark*}[Standard quantified statement]
For every $(a,b)\in\operatorname{Pre}\mathbb{Q}$,
\[
[a,b]\in\mathbb{Q}_{>0}
\Longleftrightarrow
a\cdot b>0_{\mathbb{Z}}.
\]
\end{remark*}

\begin{remark*}[Interpretation]
A rational class is declared positive when the product of numerator and
denominator in any representative is a positive integer.
\end{remark*}

\begin{dependencies}
\begin{itemize}
  \item \hyperref[def:rational-class-notation]{Rational Class Notation}
  \item \hyperref[def:multiplication-on-z-classes]{Multiplication on $\mathbb{Z}$}
  \item \hyperref[def:order-on-z]{Order on $\mathbb{Z}$}
\end{itemize}
\end{dependencies}

\begin{lemmabox}{Lemma}
\begin{lemma}[Positivity Respects Fraction Equivalence on $\mathbb{Q}$]
\label{lem:positivity-respects-equivalence-q}
If $[a,b]=[c,d]$, then
\[
a\cdot b>0_{\mathbb{Z}}
\Longleftrightarrow
c\cdot d>0_{\mathbb{Z}}.
\]
\hyperref[prf:positivity-respects-equivalence-q]{\textit{Go to proof.}}
\end{lemma}
\end{lemmabox}

\begin{remark*}[Standard quantified statement]
For all prefractions $(a,b),(c,d)$,
\[
[a,b]=[c,d]\Longrightarrow
\bigl(a\cdot b>0_{\mathbb{Z}}\Longleftrightarrow c\cdot d>0_{\mathbb{Z}}\bigr).
\]
\end{remark*}

\begin{remark*}[Interpretation]
Positivity does not depend on the chosen representative of a rational class.
\end{remark*}

\begin{dependencies}
\begin{itemize}
  \item \hyperref[thm:z-totally-ordered-set]{$\mathbb{Z}$ Is a Totally Ordered Set}
  \item \hyperref[thm:z-has-no-zero-divisors]{$\mathbb{Z}$ Has No Zero Divisors}
\end{itemize}
\end{dependencies}

\begin{theorembox}{Theorem}
\begin{theorem}[Positivity on $\mathbb{Q}$ Is Well-Defined]
\label{thm:positivity-on-q-well-defined}
The predicate ``$x$ is positive'' is well-defined on equivalence classes in
$\mathbb{Q}$.
\hyperref[prf:positivity-on-q-well-defined]{\textit{Go to proof.}}
\end{theorem}
\end{theorembox}

\begin{remark*}[Standard quantified statement]
\[
\operatorname{WellDefinedPredicate}
\bigl(\mathbb{Q}_{>0},
\operatorname{Quotient}(\operatorname{Pre}\mathbb{Q},\sim_{\mathbb{Q}})\bigr).
\]
\end{remark*}

\begin{remark*}[Interpretation]
The representative positivity test descends to a truth-valued predicate on
rational equivalence classes.
\end{remark*}

\begin{dependencies}
\begin{itemize}
  \item \hyperref[def:positive-rational]{Positive Rational}
  \item \hyperref[lem:positivity-respects-equivalence-q]{Positivity Respec
ts Fraction Equivalence on $\mathbb{Q}$}
\end{itemize}
\end{dependencies}

\begin{definitionbox}{Definition}
\begin{definition}[Strict Order on $\mathbb{Q}$]
\label{def:strict-order-on-q}
For $x,y\in\mathbb{Q}$, define
\[
x<y \quad:\Longleftrightarrow\quad y-x \text{ is positive}.
\]
\end{definition}
\end{definitionbox}

\begin{remark*}[Standard quantified statement]
For all $x,y\in\mathbb{Q}$,
\[
x<y \Longleftrightarrow y-x\in\mathbb{Q}_{>0}.
\]
\end{remark*}

\begin{remark*}[Interpretation]
Strict rational order is defined by positivity of the difference $y-x$.
\end{remark*}

\begin{dependencies}
\begin{itemize}
  \item \hyperref[def:positive-rational]{Positive Rational}
  \item \hyperref[def:subtraction-on-q]{Subtraction on $\mathbb{Q}$}
\end{itemize}
\end{dependencies}

\begin{definitionbox}{Definition}
\begin{definition}[Order on $\mathbb{Q}$]
\label{def:order-on-q}
For $x,y\in\mathbb{Q}$, define
\[
x\le y \quad:\Longleftrightarrow\quad x<y \text{ or } x=y.
\]
\end{definition}
\end{definitionbox}

\begin{remark*}[Standard quantified statement]
For all $x,y\in\mathbb{Q}$,
\[
x\le y \Longleftrightarrow x<y\lor x=y.
\]
\end{remark*}

\begin{remark*}[Interpretation]
The non-strict rational order is the reflexive closure of the strict order.
\end{remark*}

\begin{dependencies}
\begin{itemize}
  \item \hyperref[def:strict-order-on-q]{Strict Order on $\mathbb{Q}$}
\end{itemize}
\end{dependencies}

\begin{propositionbox}{Proposition}
\begin{proposition}[Product of Positives Is P
ositive on $\mathbb{Q}$]
\label{prop:product-of-positives-is-positive-q}
For all $x,y\in\mathbb{Q}$,
\[
0_{\mathbb{Q}}<x\land 0_{\mathbb{Q}}<y
\Longrightarrow
0_{\mathbb{Q}}<x\cdot y.
\]
\hyperref[prf:product-of-positives-is-positive-q]{\textit{Go to proof.}}
\end{proposition}
\end{propositionbox}

\begin{remark*}[Standard quantified statement]
For all $x,y\in\mathbb{Q}$,
\[
0_{\mathbb{Q}}<x\land 0_{\mathbb{Q}}<y
\Longrightarrow
0_{\mathbb{Q}}<x\cdot y.
\]
\end{remark*}

\begin{remark*}[Interpretation]
The product of two positive rational numbers is positive.
\end{remark*}

\begin{dependencies}
\begin{itemize}
  \item \hyperref[def:positive-rational]{Positive Rational}
  \item \hyperref[def:multiplication-on-q-classes]{Multiplication on $\mathbb{Q}$}
  \item \hyperref[thm:z-has-no-zero-divisors]{$\mathbb{Z}$ Has No Zero Divisors}
\end{itemize}
\end{dependencies}

\subsection*{\texorpdfstring{Trichotomy on $\mathbb{Q}$}{Trichotomy on Q}}

\begin{theorembox}{Theorem}
\begin{theorem}[Sign Trichotomy on $\mathbb{Q}$]
\label{thm:sign-trichotomy-on-q}
For every $x\in\mathbb{Q}$, exactly one of
\[
x<0_{\mathbb{Q}},\qquad x=0_{\mathbb{Q}},\qquad 0_{\mathbb{Q}}<x
\]
holds.
\hyperref[prf:sign-trichotomy-on-q]{\textit{Go to proof.}}
\end{theorem}
\end{theorembox}

\begin{remark*}[Standard quantified statement]
For every $x\in\mathbb{Q}$, exactly one of
\[
x<0_{\mathbb{Q}},\qquad x=0_{\mathbb{Q}},\qquad 0_{\mathbb{Q}}<x
\]
holds.
\end{remark*}

\begin{remark*}[Interpretation]
Every rational number has exactly one sign: negative, zero, or positive.
\end{remark*}

\begin{dependencies}
\begin{itemize}
  \item \hyperref[def:positive-rational]{Positive Rational}
  \item \hyperref[def:strict-order-on-q]{Strict Order on $\mathbb{Q}$}
  \item \hyperref[def:zero-in-q]{Zero in $\mathbb{Q}$}
\end{itemize}
\end{dependencies}

\begin{theorembox}{Theorem}
\begin{theorem}[Trichotomy on $\mathbb{Q}$]
\label{thm:trichotomy-on-q}
For all $x,y\in\mathbb{Q}$, exactly one of
\[
x<y,\qquad x=y,\qquad y<x
\]
holds.
\hyperref[prf:trichotomy-on-q]{\textit{Go to proof.}}
\end{theorem}
\end{theorembox}

\begin{remark*}[Standard quantified statement]
For all $x,y\in\mathbb{Q}$, exactly one of
\[
x<y,\qquad x=y,\qquad y<x
\]
holds.
\end{remark*}

\begin{remark*}[Interpretation]
Any two rational numbers are related in exactly one strict-order direction, or
they are equal.
\end{remark*}

\begin{dependencies}
\begin{itemize}
  \item \hyperref[def:strict-order-on-q]{Strict Order on $\mathbb{Q}$}
  \item \hyperref[def:subtraction-on-q]{Subtraction on $\mathbb{Q}$}
  \item \hyperref[thm:sign-trichotomy-on-q]{Sign Trichotomy on $\mathbb{Q}$}
\end{itemize}
\end{dependencies}

\subsection*{\texorpdfstring{Total Order on $\mathbb{Q}$}{Total Order on Q}}

\begin{theorembox}{Theorem}
\begin{theorem}[Strict Order on $\mathbb{Q}$ Is Irreflexive]
\label{thm:strict-order-on-q-irreflexive}
For every $x\in\mathbb{Q}$, not $x<x$.
\hyperref[prf:strict-order-on-q-irreflexive]{\textit{Go to proof.}}
\end{theorem}
\end{theorembox}

\begin{remark*}[Standard quantified statement]
For every $x\in\mathbb{Q}$,
\[
\neg(x<x).
\]
\end{remark*}

\begin{remark*}[Interpretation]
No rational number is strictly less than itself.
\end{remark*}

\begin{dependencies}
\begin{itemize}
  \item \hyperref[def:strict-order-on-q]{Strict Order on $\mathbb{Q}$}
  \item \hyperref[thm:trichotomy-on-q]{Trichotomy on $\mathbb{Q}$}
\end{itemize}
\end{dependencies}

\begin{theorembox}{Theorem}
\begin{theorem}[Strict Order on $\mathbb{Q}$ Is Asymmetric]
\label{thm:strict-order-on-q-asymmetric}
For all $x,y\in\mathbb{Q}$,
\[
x<y \Longrightarrow \neg(y<x).
\]
\hyperref[prf:strict-order-on-q-asymmetric]{\textit{Go to proof.}}
\end{theorem}
\end{theorembox}

\begin{remark*}[Standard quantified statement]
For all $x,y\in\mathbb{Q}$,
\[
x<y\Longrightarrow \neg(y<x).
\]
\end{remark*}

\begin{remark*}[Interpretation]
Strict rational order cannot hold in both directions between the same two
numbers.
\end{remark*}

\begin{dependencies}
\begin{itemize}
  \item \hyperref[def:strict-order-on-q]{Strict Order on $\mathbb{Q}$}
  \item \hyperref[thm:trichotomy-on-q]{Trichotomy on $\mathbb{Q}$}
\end{itemize}
\end{dependencies}

\begin{theorembox}{Theorem}
\begin{theorem}[Strict Order on $\mathbb{Q}$ Is Transitive]
\label{thm:strict-order-on-q-transitive}
For all $x,y,z\in\mathbb{Q}$,
\[
x
<y\land y<z \Longrightarrow x<z.
\]
\hyperref[prf:strict-order-on-q-transitive]{\textit{Go to proof.}}
\end{theorem}
\end{theorembox}

\begin{remark*}[Standard quantified statement]
For all $x,y,z\in\mathbb{Q}$,
\[
x<y\land y<z\Longrightarrow x<z.
\]
\end{remark*}

\begin{remark*}[Interpretation]
Strict rational order is transitive.
\end{remark*}

\begin{dependencies}
\begin{itemize}
  \item \hyperref[def:strict-order-on-q]{Strict Order on $\mathbb{Q}$}
  \item \hyperref[thm:addition-preserves-strict-order-on-q]{Addition Preserves Strict Order on $\mathbb{Q}$}
\end{itemize}
\end{dependencies}

\begin{theorembox}{Theorem}
\begin{theorem}[Order on $\mathbb{Q}$ Is Reflexive]
\label{thm:order-on-q-reflexive}
For every $x\in\mathbb{Q}$, $x\le x$.
\hyperref[prf:order-on-q-reflexive]{\textit{Go to proof.}}
\end{theorem}
\end{theorembox}

\begin{remark*}[Standard quantified statement]
For every $x\in\mathbb{Q}$,
\[
x\le x.
\]
\end{remark*}

\begin{remark*}[Interpretation]
The non-strict rational order is reflexive.
\end{remark*}

\begin{dependencies}
\begin{itemize}
  \item \hyperref[def:order-on-q]{Order on $\mathbb{Q}$}
\end{itemize}
\end{dependencies}

\begin{theorembox}{Theorem}
\begin{theorem}[Order on $\mathbb{Q}$ Is Antisymmetric]
\label{thm:order-on-q-antisymmetric}
For all $x,y\in\mathbb{Q}$,
\[
x\le y\land y\le x \Longrightarrow x=y.
\]
\hyperref[prf:order-on-q-antisymmetric]{\textit{Go to proof.}}
\end{theorem}
\end{theorembox}

\begin{remark*}[Standard quantified statement]
For all $x,y\in\mathbb{Q}$,
\[
x\le y\land y\le x\Longrightarrow x=y.
\]
\end{remark*}

\begin{remark*}[Interpretation]
The non-strict rational order is antisymmetric.
\end{remark*}

\begin{dependencies}
\begin{itemize}
  \item \hyperref[def:order-on-q]{Order on $\mathbb{Q}$}
  \item \hyperref[thm:strict-order-on-q-asymmetric]{Strict Order on $\mathbb{Q}$ Is Asymmetric}
\end{itemize}
\end{dependencies}

\begin{theorembox}{Theorem}
\begin{theorem}[Order on $\mathbb{Q}$ Is Transitive]
\label{thm:order-on-q-transitive}
For all $x,y,z\in\mathbb{Q}$,
\[
x\le y\land y\le z \Longrightarrow x\le z.
\]
\hyperref[prf:order-on-q-transitive]{\textit{Go to proof.}}
\end{theorem}
\end{theorembox}

\begin{remark*}[Standard quantified statement]
For all $x,y,z\in\mathbb{Q}$,
\[
x\le y\land y\le z\Longrightarrow x\le z.
\]
\end{remark*}

\begin{remark*}[Interpretation]
The non-strict rational order is transitive.
\end{remark*}

\begin{dependencies}
\begin{itemize}
  \item \hyperref[def:order-on-q]{Order on $\mathbb{Q}$}
  \item \hyperref[thm:strict-order-on-q-transitive]{Strict Order on $\mathbb{Q}$ Is Transitive}
\end{itemize}
\end{dependencies}

\begin{theorembox}{Theorem}
\begin{theorem}[Order on $\mathbb{Q}$ Is Total]
\label{thm:order-on-q-total}
For all $x,y\in\mathbb{Q}$,
\[
x\le y\text{ or }y\le x.
\]
\hyperref[prf:order-on-q-total]{\textit{Go to proof.}}
\end{theorem}
\end{theorembox}

\begin{remark*}[Standard quantified statement]
For all $x,y\in\mathbb{Q}$,
\[
x\le y\lor y\le x.
\]
\end{remark*}

\begin{remark*}[Interpretation]
The non-strict rational order compares every pair of rational numbers.
\end{remark*}

\begin{dependencies}
\begin{itemize}
  \item \hyperref[def:order-on-q]{Order on $\mathbb{Q}$}
  \item \hyperref[thm:trichotomy-on-q]{Trichotomy on $\mathbb{Q}$}
\end{itemize}
\end{dependencies}

\begin{theorembox}{Theorem}
\begin{theorem}[$\mathbb{Q}$ Is a Totally Ordered Set]
\label{thm:q-totally-ordered-set}
The ordered pair $(\mathbb{Q},\le)$ is a totally ordered set.
\hyperref[prf:q-totally-ordered-set]{\textit{Go to proof.}}
\end{theorem}
\end{theorembox}

\begin{remark*}[Standard quantified statement]
\[
\operatorname{TotalOrder}(\mathbb{Q},\le).
\]
\end{remark*}

\begin{remark*}[Interpretation]
The rational carrier together with its non-strict order forms a totally
ordered set.
\end{remark*}

\begin{dependencies}
\begin{itemize}
  \item \hyperref[def:ordered-set]{Ordered Set}
  \item \hyperref[def:total-order]{Total Order}
  \item \hyperref[def:order-on-q]{Order on $\mathbb{Q}$}
  \item \hyperref[thm:order-on-q-reflexive]{Order on $\mathbb{Q}$ Is Reflexive}
  \item \hyperref[thm:order-on-q-antisymmetric]{Order on $\mathbb{Q}$ Is Antisymmetric}
  \item \hyperref[thm:order-on-q-transitive]{Order on $\mathbb{Q}$ Is Transitive}
  \item \hyperref[thm:order-on-q-total]{Order on $\mathbb{Q}$ Is Total}
\end{itemize}
\end{dependencies}

\subsection*{\texorpdfstring{Addition Order Compatibility on $\mathbb{Q}$}{Addition Order Compatibility on Q}}

\begin{theorembox}{Theorem}
\begin{theorem}[Left Addition Preserves Strict Order on $\mathbb{Q}$]
\label{thm:left-addition-preserves-strict-order-on-q}
For all $x,y,z\in\mathbb{Q}$,
\[
x<y \Longrightarrow z+x<z+y.
\]
\hyperref[prf:left-addition-preserves-strict-order-on-q]{\textit{Go to proof.}}
\end{theorem}
\end{theorembox}

\begin{remark*}[Standard quantified statement]
For all $x,y,z\in\mathbb{Q}$,
\[
x<y\Longrightarrow z+x<z+y.
\]
\end{remark*}

\begin{remark*}[Interpretation]
Left addition by a fixed rational preserves strict order.
\end{remark*}

\begin{dependencies}
\begin{itemize}
  \item \hyperref[def:strict-order-on-q]{Strict Order on $\mathbb{Q}$}
  \item \hyperref[def:addition-on-q-classes]{Addition on $\mathbb{Q}$}
  \item \hyperref[def:positive-rational]{Positive Rational}
\end{itemize}
\end{dependencies}

\begin{corollarybox}{Corollary}
\begin{corollary}[Right Addition Preserves Strict Order on $\mathbb{Q}$]
\label{cor:right-addition-preserves-strict-order-on-q}
For all $x,y,z\in\mathbb{Q}$,
\[
x<y \Longrightarrow x+z<y+z.
\]
\hyperref[prf:right-addition-preserves-strict-order-on-q]{\textit{Go to proof.}}
\end{corollary}
\end{corollarybox}

\begin{remark*}[Standard quantified statement]
For all $x,y,z\in\mathbb{Q}$,
\[
x<y\Longrightarrow x+z<y+z.
\]
\end{remark*}

\begin{remark*}[Interpretation]
Right addition by a fixed rational preserves strict order.
\end{remark*}

\begin{dependencies}
\begin{itemize}
  \item \hyperref[thm:left-addition-preserves-strict-order-on-q]{Left Addition Preserves Strict Order on $\mathbb{Q}$}
  \item \hyperref[thm:addition-on-q-commutative]{Addition on $\mathbb{Q}$ Is Commutative}
\end{itemize}
\end{dependencies}

\begin{theorembox}{Theorem}
\begin{theorem}[Addition Preserves Strict Order on $\mathbb{Q}$]
\label{thm:addition-preserves-strict-order-on-q}
Addition preserves strict order on both sides in $\mathbb{Q}$.
\hyperref[prf:addition-preserves-strict-order-on-q]{\textit{Go to proof.}}
\end{theorem}
\end{theorembox}

\begin{remark*}[Standard quantified statement]
For all $x,y,z\in\mathbb{Q}$,
\[
x<y\Longrightarrow z+x<z+y
\quad\text{and}\quad
x<y\Longrightarrow x+z<y+z.
\]
\end{remark*}

\begin{remark*}[Interpretation]
Addition by a fixed rational preserves strict order on both sides.
\end{remark*}

\begin{dependencies}
\begin{itemize}
  \item \hyperref[thm:left-addition-preserves-strict-order-on-q]{Left Addition Preserves Strict Order on $\mathbb{Q}$}
  \item \hyperref[cor:right-addition-preserves-strict-order-on-q]{Right Addition Preserves Strict Order on $\mathbb{Q}$}
\end{itemize}
\end{dependencies}

\begin{theorembox}{Theorem}
\begin{theorem}[Left Addition Preserves Order on $\mathbb{Q}$]
\label{thm:left-addition-preserves-order-on-q}
For all $x,y,z\in\mathbb{Q}$,
\[
x\le y \Longrightarrow z+x\le z+y.
\]
\hyperref[prf:left-addition-preserves-order-on-q]{\textit{Go to proof.}}
\end{theorem}
\end{theorembox}

\begin{remark*}[Standard quantified statement]
For all $x,y,z\in\mathbb{Q}$,
\[
x\le y\Longrightarrow z+x\le z+y.
\]
\end{remark*}

\begin{remark*}[Interpretation]
Left addition by a fixed rational preserves non-strict order.
\end{remark*}

\begin{dependencies}
\begin{itemize}
  \item \hyperref[def:order-on-q]{Order on $\mathbb{Q}$}
  \item \hyperref[thm:left-addition-preserves-strict-order-on-q]{Left Addition Preserves Strict Order on $\mathbb{Q}$}
\end{itemize}
\end{dependencies}

\begin{corollarybox}{Corollary}
\begin{corollary}[Right Addition Preserves Order on $\mathbb{Q}$]
\label{cor:right-addition-preserves-order-on-q}
For all $x,y,z\in\mathbb{Q}$,
\[
x\le y \Longrightarrow x+z\le y+z.
\]
\hyperref[prf:right-addition-preserves-order-on-q]{\textit{Go to proof.}}
\end{corollary}
\end{corollarybox}

\begin{remark*}[Standard quantified statement]
For all $x,y,z\in\mathbb{Q}$,
\[
x\le y\Longrightarrow x+z\le y+z.
\]
\end{remark*}

\begin{remark*}[Interpretation]
Right addition by a fixed rational preserves non-strict order.
\end{remark*}

\begin{dependencies}

\begin{itemize}
  \item \hyperref[def:order-on-q]{Order on $\mathbb{Q}$}
  \item \hyperref[cor:right-addition-preserves-strict-order-on-q]{Right Addition Preserves Strict Order on $\mathbb{Q}$}
\end{itemize}
\end{dependencies}

\begin{theorembox}{Theorem}
\begin{theorem}[Addition Preserves Order on $\mathbb{Q}$]
\label{thm:addition-preserves-order-on-q}
Addition preserves non-strict order on both sides in $\mathbb{Q}$.
\hyperref[prf:addition-preserves-order-on-q]{\textit{Go to proof.}}
\end{theorem}
\end{theorembox}

\begin{remark*}[Standard quantified statement]
For all $x,y,z\in\mathbb{Q}$,
\[
x\le y\Longrightarrow z+x\le z+y
\quad\text{and}\quad
x\le y\Longrightarrow x+z\le y+z.
\]
\end{remark*}

\begin{remark*}[Interpretation]
Addition by a fixed rational preserves non-strict order on both sides.
\end{remark*}

\begin{dependencies}
\begin{itemize}
  \item \hyperref[thm:left-addition-preserves-order-on-q]{Left Addition Preserves Order on $\mathbb{Q}$}
  \item \hyperref[cor:right-addition-preserves-order-on-q]{Right Addition Preserves Order on $\mathbb{Q}$}
\end{itemize}
\end{dependencies}

\begin{corollarybox}{Corollary}
\begin{corollary}[Addition Reflects Order on $\mathbb{Q}$]
\label{cor:addition-reflects-order-on-q}
Addition by a fixed rational reflects
 both strict and non-strict order.
\hyperref[prf:addition-reflects-order-on-q]{\textit{Go to proof.}}
\end{corollary}
\end{corollarybox}

\begin{remark*}[Standard quantified statement]
For all $x,y,z\in\mathbb{Q}$,
\[
z+x<z+y\Longrightarrow x<y,
\qquad
z+x\le z+y\Longrightarrow x\le y.
\]
\end{remark*}

\begin{remark*}[Interpretation]
Translation by a fixed rational is order-reflecting because it can be undone
by adding the additive inverse.
\end{remark*}

\begin{dependencies}
\begin{itemize}
  \item \hyperref[thm:addition-preserves-strict-order-on-q]{Addition Preserves Strict Order on $\mathbb{Q}$}
  \item \hyperref[thm:addition-preserves-order-on-q]{Addition Preserves Order on $\mathbb{Q}$}
  \item \hyperref[thm:negation-left-additive-inverse-on-q]{Negation Gives a Left Additive Inverse on $\mathbb{Q}$}
\end{itemize}
\end{dependencies}

\subsection*{\texorpdfstring{Multiplication Order Compatibility on $\mathbb{Q}$}{Multiplication Order Compatibility on Q}}

\begin{theorembox}{Theorem}
\begin{theorem}[Left Multiplication by Positives Preserve
s Strict Order on $\mathbb{Q}$]
\label{thm:left-positive-multiplication-preserves-strict-order-on-q}
For all $x,y,z\in\mathbb{Q}$,
\[
0_{\mathbb{Q}}<z\land x<y
\Longrightarrow
z\cdot x<z\cdot y.
\]
\hyperref[prf:left-positive-multiplication-preserves-strict-order-on-q]{\textit{Go to proof.}}
\end{theorem}
\end{theorembox}

\begin{remark*}[Standard quantified statement]
For all $x,y,z\in\mathbb{Q}$,
\[
0_{\mathbb{Q}}<z\land x<y
\Longrightarrow
z\cdot x<z\cdot y.
\]
\end{remark*}

\begin{remark*}[Interpretation]
Left multiplication by a positive rational preserves strict order.
\end{remark*}

\begin{dependencies}
\begin{itemize}
  \item \hyperref[def:strict-order-on-q]{Strict Order on $\mathbb{Q}$}
  \item \hyperref[prop:product-of-positives-is-positive-q]{Product of Positives Is Positive on $\mathbb{Q}$}
  \item \hyperref[def:multiplication-on-q-classes]{Multiplication on $\mathbb{Q}$}
\end{itemize}
\end{dependencies}

\begin{corollarybox}{Corollary}
\begin{corollary}[Right Multiplication by Positives Preserves Strict Order on $\mathbb{Q}$]
\label{cor:right-positive-multiplication-preserves-strict-order-on-q}
For all $x,y,z\in\mathbb{Q}$,
\[
0_{\mathbb{Q}}<z\land x<y
\Longrightarrow
x\cdot z<y\cdot z.
\]
\hyperref[prf:right-positive-multiplication-preserves-strict-order-on-q]{\textit{Go to proof.}}
\end{corollary}
\end{corollarybox}

\begin{remark*}[Standard quantified statement]
For all $x,y,z\in\mathbb{Q}$,
\[
0_{\mathbb{Q}}<z\land x<y
\Longrightarrow
x\cdot z<y\cdot z.
\]
\end{remark*}

\begin{remark*}[Interpretation]
Right multiplication by a positive rational preserves strict order.
\end{remark*}

\begin{dependencies}
\begin{itemize}
  \item \hyperref[thm:left-positive-multiplication-preserves-strict-order-on-q]{Left Multiplication by Positives Preserves Strict Order on $\mathbb{Q}$}
  \item \hyperref[thm:multiplication-on-q-commutative]{Multiplication on $\mathbb{Q}$ Is Commutative}
\end{itemize}
\end{dependencies}

\begin{theorembox}{Theorem}
\begin{theorem}[Multiplication by Positives Preserves Strict Order on $\mathbb{Q}$]
\label{thm:positive-multiplication-preserves-strict-order-on-q}
Multiplication by a positive rational preserves strict order on both sides.
\hyperref[prf:positive-multiplication-preserves-strict-order-on-q]{\textit{Go to proof.}}
\end{theorem}
\end{theorembox}

\begin{remark*}[Standard quantified statement]
For all $x,y,z\in\mathbb{Q}$,
\[
0_{\mathbb{Q}}<z\land x<y
\Longrightarrow
z\cdot x<z\cdot y
\quad\text{and}\quad
x\cdot z<y\cdot z.
\]
\end{remark*}

\begin{remark*}[Interpretation]
Multiplication by a positive rational preserves strict order on both sides.
\end{remark*}

\begin{dependencies}
\begin{itemize}
  \item \hyperref[thm:left-positive-multiplication-preserves-strict-order-on-q]{Left Multiplication by Positives Preserves Strict Order on $\mathbb{Q}$}
  \item \hyperref[cor:right-positive-multiplication-preserves-strict-order-on-q]{Right Multiplication by Positives Preserves Strict Order on $\mathbb{Q}$}
\end{itemize}
\end{dependencies}

\begin{theorembox}{Theorem}
\begin{theorem}[Left Multiplication by Positives Preserves Order on $\mathbb{Q}$]
\label{thm:left-positive-multiplication-preserves-order-on-q}
For all $x,y,z\in\mathbb{Q}$,
\[
0_{\mathbb{Q}}<z\land x\le y
\Longrightarrow
z\cdot x\le z\cdot y.
\]
\hyperref[prf:left-positive-multiplication-preserves-order-on-q]{\textit{Go to proof.}}
\end{theorem}
\end{theorembox}

\begin{remark*}[Standard quantified statement]
For all $x,y,z\in\mathbb{Q}$,
\[
0_{\mathbb{Q}}<z\land x\le y
\Longrightarrow
z\cdot x\le z\cdot y.
\]
\end{remark*}

\begin{remark*}[Interpretation]
Left multiplication by a positive rational preserves non-strict order.
\end{remark*}

\begin{dependencies}
\begin{itemize}
  \item \hyperref[def:order-on-q]{Order on $\mathbb{Q}$}
  \item \hyperref[thm:left-positive-multiplication-preserves-strict-order-on-q]{Left Multiplication by Positives Preserves Strict Order on $\mathbb{Q}$}
\end{itemize}
\end{dependencies}

\begin{corollarybox}{Corollary}
\begin{corollary}[Right Multiplication by Positives Preserves Order
 on $\mathbb{Q}$]
\label{cor:right-positive-multiplication-preserves-order-on-q}
For all $x,y,z\in\mathbb{Q}$,
\[
0_{\mathbb{Q}}<z\land x\le y
\Longrightarrow
x\cdot z\le y\cdot z.
\]
\hyperref[prf:right-positive-multiplication-preserves-order-on-q]{\textit{Go to proof.}}
\end{corollary}
\end{corollarybox}

\begin{remark*}[Standard quantified statement]
For all $x,y,z\in\mathbb{Q}$,
\[
0_{\mathbb{Q}}<z\land x\le y
\Longrightarrow
x\cdot z\le y\cdot z.
\]
\end{remark*}

\begin{remark*}[Interpretation]
Right multiplication by a positive rational preserves non-strict order.
\end{remark*}

\begin{dependencies}
\begin{itemize}
  \item \hyperref[def:order-on-q]{Order on $\mathbb{Q}$}
  \item \hyperref[cor:right-positive-multiplication-preserves-strict-order-on-q]{Right Multiplication by Positives Preserves Strict Order on $\mathbb{Q}$}
\end{itemize}
\end{dependencies}

\begin{theorembox}{Theorem}
\begin{theorem}[Multiplication by Positives Preserves Order on $\mathbb{Q}$]
\label{thm:positive-multiplication-preserves-order-on-q}
Multiplication by a positive rational preserves non-strict order on both sides.
\hyperref[prf:positive-multiplication-preserves-order-on-q]{\textit{Go to proof.}}
\end{theorem}
\end{theorembox}

\begin{remark*}[Standard quantified statement]
For all $x,y,z\in\mathbb{Q}$,
\[
0_{\mathbb{Q}}<z\land x\le y
\Longrightarrow
z\cdot x\le z\cdot y
\quad\text{and}\quad
x\cdot z\le y\cdot z.
\]
\end{remark*}

\begin{remark*}[Interpretation]
Multiplication by a positive rational preserves non-strict order on both
sides.
\end{remark*}

\begin{dependencies}
\begin{itemize}
  \item \hyperref[thm:left-positive-multiplication-preserves-order-on-q]{Left Multiplication by Positives Preserves Order on $\mathbb{Q}$}
  \item \hyperref[cor:right-positive-multiplication-preserves-order-on-q]{Right Multiplication by Positives Preserves Order on $\mathbb{Q}$}
\end{itemize}
\end{dependencies}

\begin{theorembox}{Theorem}
\begin{theorem}[Multiplication by Negatives Reverses Strict Order on $\mathbb{Q}$]
\label{thm:negative-multiplication-reverses-strict-order-on-q}
For all $x,y,z\in\mathbb{Q}$,
\[
z<0_{\mathbb{Q}}\land x<y
\Longrightarrow
z\cdot y<z\cdot x.
\]
\hyperref[prf:negative-multiplication-reverses-strict-order-on-q]{\textit{Go to proof.}}
\end{theorem}
\end{theorembox}

\begin{remark*}[Standard quantified statement]
For all $x,y,z\in\mathbb{Q}$,
\[
z<0_{\mathbb{Q}}\land x<y
\Longrightarrow
z\cdot y<z\cdot x.
\]
\end{remark*}

\begin{remark*}[Interpretation]
Left multiplication by a negative rational reverses strict order.
\end{remark*}

\begin{dependencies}
\begin{itemize}
  \item \hyperref[def:strict-order-on-q]{Strict Order on $\mathbb{Q}$}
  \item \hyperref[cor:negation-of-product-on-q]{Negation of a Product on $\mathbb{Q}$}
  \item \hyperref[thm:positive-multiplication-preserves-strict-order-on-q]{Positive Multiplication Preserves Strict Order on $\mathbb{Q}$}
\end{itemize}
\end{dependencies}

\begin{corollarybox}{Corollary}
\begin{corollary}[Multiplication by Negatives Reverses Order on $\mathbb{Q}$]
\label{cor:negative-multiplication-reverses-order-on-q}
For all $x,y,z\in\mathbb{Q}$,
\[
z<0_{\mathbb{Q}}\land x\le y
\Longrightarrow
z\cdot y\le z\cdot x.
\]
\hyperref[prf:negative-multiplication-reverses-order-on-q]{\textit{Go to proof.}}
\end{corollary}
\end{corollarybox}

\begin{remark*}[Standard quantified statement]
For all $x,y,z\in\mathbb{Q}$,
\[
z<0_{\mathbb{Q}}\land x\le y
\Longrightarrow
z\cdot y\le z\cdot x.
\]
\end{remark*}

\begin{remark*}[Interpretation]
Left multiplication by a negative rational reverses non-strict order.
\end{remark*}

\begin{dependencies}
\begin{itemize}
  \item \hyperref[def:order-on-q]{Order on $\mathbb{Q}$}
  \item \hyperref[thm:negative-multiplication-reverses-strict-order-on-q]{Multiplication by Negatives Reverses Strict Order on $\mathbb{Q}$}
\end{itemize}
\end{dependencies}

\subsection*{\texorpdfstring{Reciprocal Order Laws on $\mathbb{Q}$}{Reciprocal Order Laws on Q}}

\begin{theorembox}{Theorem}
\begin{theorem}[Sign of a Reciprocal on $\mathbb{Q}$]
\label{thm:sign-of-reciprocal-q}
For every nonzero $x\in\mathbb{Q}$,
\[
0_{\mathbb{Q}}<x \Longrightarrow 0_{\mathbb{Q}}<x^{-1},
\qquad
x<0_{\mathbb{Q}} \Longrightarrow x^{-1}<0_{\mathbb{Q}}.
\]
\hyperref[prf:sign-of-reciprocal-q]{\textit{Go to proof.}}
\end{theorem}
\end{theorembox}

\begin{remark*}[Standard quantified statement]
For every $x\in\mathbb{Q}\setminus\{0_{\mathbb{Q}}\}$,
\[
0_{\mathbb{Q}}<x\Longrightarrow 0_{\mathbb{Q}}<x^{-1},
\qquad
x<0_{\mathbb{Q}}\Longrightarrow x^{-1}<0_{\mathbb{Q}}.
\]
\end{remark*}

\begin{remark*}[Interpretation]
A nonzero rational number and its reciprocal have the same sign.
\end{remark*}

\begin{dependencies}
\begin{itemize}
  \item \hyperref[def:reciprocal-on-q-classes]{Reciprocal on $\mathbb{Q}$}
  \item \hyperref[def:order-on-q]{Order on $\mathbb{Q}$}
  \item \hyperref[prop:product-of-positives-is-positive-q]{Product of Positives Is Positive on $\mathbb{Q}$}
  \item \hyperref[thm:reciprocal-left-multiplicative-inverse-on-q]{Reciprocal Is a Left Multiplicative Inverse on $\mathbb{Q}$}
\end{itemize}
\end{dependencies}

\begin{theorembox}{Theorem}
\begin{theorem}[Reciprocal Reverses Order on Positives in $\mathbb{Q}$]
\label{thm:reciprocal-reverses-order-on-positives-q}
For all $x,y\in\mathbb{Q}$,
\[
0_{\mathbb{Q}}<x<y
\Longrightarrow
0_{\mathbb{Q}}<y^{-1}<x^{-1}.
\]
\hyperref[prf:reciprocal-reverses-order-on-positives-q]{\textit{Go to proof.}}
\end{theorem}
\end{theorembox}

\begin{remark*}[Standard quantified statement]
For all $x,y\in\mathbb{Q}$,
\[
0_{\mathbb{Q}}<x<y
\Longrightarrow
0_{\mathbb{Q}}<y^{-1}<x^{-1}.
\]
\end{remark*}

\begin{remark*}[Interpretation]
On positive rationals, taking reciprocals reverses strict order.
\end{remark*}

\begin{dependencies}
\begin{itemize}
  \item \hyperref[def:reciprocal-on-q-classes]{Reciprocal on $\mathbb{Q}$}
  \item \hyperref[def:order-on-q]{Order on $\mathbb{Q}$}
  \item \hyperref[thm:sign-of-reciprocal-q]{Sign of a Reciprocal on $\mathbb{Q}$}
  \item \hyperref[thm:positive-multiplication-preserves-strict-order-on-q]{Positive Multiplication Preserves Strict Order on $\mathbb{Q}$}
\end{itemize}
\end{dependencies}

\begin{theorembox}{Theorem}
\begin{theorem}[Reciprocal Reverses Order on Same-Sign Rationals]
\label{thm:reciprocal-reverses-order-on-same-sign-q}
For all nonzero $x,y\in\mathbb{Q}$,
\[
0_{\mathbb{Q}}<x<y
\Longrightarrow
0_{\mathbb{Q}}<y^{-1}<x^{-1},
\qquad
x<y<0_{\mathbb{Q}}
\Longrightarrow
y^{-1}<x^{-1}<0_{\mathbb{Q}}.
\]
\hyperref[prf:reciprocal-reverses-order-on-same-sign-q]{\textit{Go to proof.}}
\end{theorem}
\end{theorembox}

\begin{remark*}[Standard quantified statement]
For all nonzero $x,y\in\mathbb{Q}$,
\[
0_{\mathbb{Q}}<x<y
\Longrightarrow
0_{\mathbb{Q}}<y^{-1}<x^{-1},
\qquad
x<y<0_{\mathbb{Q}}
\Longrightarrow
y^{-1}<x^{-1}<0_{\mathbb{Q}}.
\]
\end{remark*}

\begin{remark*}[Interpretation]
For rationals with the same sign, reciprocal reverses strict order and
preserves the common sign side.
\end{remark*}

\begin{dependencies}
\begin{itemize}
  \item \hyperref[def:reciprocal-on-q-classes]{Reciprocal on $\mathbb{Q}$}
  \item \hyperref[def:order-on-q]{Order on $\mathbb{Q}$}
  \item \hyperref[thm:sign-of-reciprocal-q]{Sign of a Reciprocal on $\mathbb{Q}$}
  \item \hyperref[thm:reciprocal-reverses-order-on-positives-q]{Reciprocal Reverses Order on Positives in $\mathbb{Q}$}
  \item \hyperref[thm:negative-multiplication-reverses-strict-order-on-q]{Multiplication by Negatives Reverses Strict Order on $\mathbb{Q}$}
\end{itemize}
\end{dependencies}

\begin{theorembox}{Theorem}
\begin{theorem}[Reciprocal Is Involutive on $\mathbb{Q}$]
\label{thm:reciprocal-involutive-on-q}
For every nonzero $x\in\mathbb{Q}$,
\[
(x^{-1})^{-1}=x.
\]
\hyperref[prf:reciprocal-involutive-on-q]{\textit{Go to proof.}}
\end{theorem}
\end{theorembox}

\begin{remark*}[Standard quantified statement]
For every $x\in\mathbb{Q}\setminus\{0_{\mathbb{Q}}\}$,
\[
(x^{-1})^{-1}=x.
\]
\end{remark*}

\begin{remark*}[Interpretation]
Taking the reciprocal twice returns the original nonzero rational number.
\end{remark*}

\begin{dependencies}
\begin{itemize}
  \item \hyperref[def:reciprocal-on-q-classes]{Reciprocal on $\mathbb{Q}$}
  \item \hyperref[thm:reciprocal-on-q-well-defined]{Reciprocal on $\mathbb{Q}$ Is Well-Defined}
  \item \hyperref[thm:reciprocal-left-multiplicative-inverse-on-q]{Reciprocal Is a Left Multiplicative Inverse on $\mathbb{Q}$}
  \item \hyperref[cor:reciprocal-right-multiplicative-inverse-on-q]{Reciprocal Is a Right Multiplicative Inverse on $\mathbb{Q}$}
\end{itemize}
\end{dependencies}

\begin{theorembox}{Theorem}
\begin{theorem}[Reciprocal of a Product on $\mathbb{Q}$]
\label{thm:reciprocal-of-product-on-q}
For all nonzero $x,y\in\mathbb{Q}$,
\[
(xy)^{-1}=x^{-1}y^{-1}.
\]
\hyperref[prf:reciprocal-of-product-on-q]{\textit{Go to proof.}}
\end{theorem}
\end{theorembox}

\begin{remark*}[Standard quantified statement]
For all $x,y\in\mathbb{Q}\setminus\{0_{\mathbb{Q}}\}$,
\[
(xy)^{-1}=x^{-1}y^{-1}.
\]
\end{remark*}

\begin{remark*}[Interpretation]
The reciprocal of a product of nonzero rational numbers is the product of the
reciprocals.
\end{remark*}

\begin{dependencies}
\begin{itemize}
  \item \hyperref[def:reciprocal-on-q-classes]{Reciprocal on $\mathbb{Q}$}
  \item \hyperref[thm:reciprocal-on-q-well-defined]{Reciprocal on $\mathbb{Q}$ Is Well-Defined}
  \item \hyperref[thm:reciprocal-left-multiplicative-inverse-on-q]{Reciprocal Is a Left Multiplicative Inverse on $\mathbb{Q}$}
  \item \hyperref[cor:reciprocal-right-multiplicative-inverse-on-q]{Reciprocal Is a Right Multiplicative Inverse on $\mathbb{Q}$}
  \item \hyperref[thm:multiplication-on-q-associative]{Multiplication on $\mathbb{Q}$ Is Associative}
  \item \hyperref[thm:multiplication-on-q-commutative]{Multiplication on $\mathbb{Q}$ Is Commutative}
\end{itemize}
\end{dependencies}

\begin{theorembox}{Theorem}
\begin{theorem}[Reciprocal Reverses Non-Strict Order on Positives in $\mathbb{Q}$]
\label{thm:reciprocal-reverses-nonstrict-order-on-positives-q}
For all nonzero $x,y\in\mathbb{Q}$,
\[
0_{\mathbb{Q}}<x\le y
\Longrightarrow
0_{\mathbb{Q}}<y^{-1}\le x^{-1}.
\]
\hyperref[prf:reciprocal-reverses-nonstrict-order-on-positives-q]{\textit{Go to proof.}}
\end{theorem}
\end{theorembox}

\begin{remark*}[Standard quantified statement]
For all $x,y\in\mathbb{Q}\setminus\{0_{\mathbb{Q}}\}$,
\[
0_{\mathbb{Q}}<x\le y
\Longrightarrow
0_{\mathbb{Q}}<y^{-1}\le x^{-1}.
\]
\end{remark*}

\begin{remark*}[Interpretation]
On positive rationals, taking reciprocals reverses non-strict order.
\end{remark*}

\begin{dependencies}
\begin{itemize}
  \item \hyperref[def:reciprocal-on-q-classes]{Reciprocal on $\mathbb{Q}$}
  \item \hyperref[def:order-on-q]{Order on $\mathbb{Q}$}
  \item \hyperref[thm:sign-of-reciprocal-q]{Sign of a Reciprocal on $\mathbb{Q}$}
  \item \hyperref[thm:reciprocal-reverses-order-on-positives-q]{Reciprocal Reverses Order on Positives in $\mathbb{Q}$}
\end{itemize}
\end{dependencies}

\begin{theorembox}{Theorem}
\begin{theorem}[Reciprocal Reverses Non-Strict Order on Same-Sign Rationals]
\label{thm:reciprocal-reverses-nonstrict-order-on-same-sign-q}
For all nonzero $x,y\in\mathbb{Q}$,
\[
0_{\mathbb{Q}}<x\le y
\Longrightarrow
0_{\mathbb{Q}}<y^{-1}\le x^{-1},
\qquad
x\le y<0_{\mathbb{Q}}
\Longrightarrow
y^{-1}\le x^{-1}<0_{\mathbb{Q}}.
\]
\hyperref[prf:reciprocal-reverses-nonstrict-order-on-same-sign-q]{\textit{Go to proof.}}
\end{theorem}
\end{theorembox}

\begin{remark*}[Standard quantified statement]
For all nonzero $x,y\in\mathbb{Q}$,
\[
0_{\mathbb{Q}}<x\le y
\Longrightarrow
0_{\mathbb{Q}}<y^{-1}\le x^{-1},
\qquad
x\le y<0_{\mathbb{Q}}
\Longrightarrow
y^{-1}\le x^{-1}<0_{\mathbb{Q}}.
\]
\end{remark*}

\begin{remark*}[Interpretation]
For rationals with the same sign, reciprocal reverses non-strict order and
preserves the common sign side.
\end{remark*}

\begin{dependencies}
\begin{itemize}
  \item \hyperref[def:reciprocal-on-q-classes]{Reciprocal on $\mathbb{Q}$}
  \item \hyperref[def:order-on-q]{Order on $\mathbb{Q}$}
  \item \hyperref[thm:sign-of-reciprocal-q]{Sign of a Reciprocal on $\mathbb{Q}$}
  \item \hyperref[thm:reciprocal-reverses-nonstrict-order-on-positives-q]{Reciprocal Reverses Non-Strict Order on Positives in $\mathbb{Q}$}
  \item \hyperref[cor:negative-multiplication-reverses-order-on-q]{Multiplication by Negatives Reverses Order on $\mathbb{Q}$}
\end{itemize}
\end{dependencies}

\subsection*{\texorpdfstring{Absolute Value on $\mathbb{Q}$}{Absolute Value on Q}}

\begin{definitionbox}{Definition}
\begin{definition}[Absolute Value on $\mathbb{Q}$
]
\label{def:absolute-value-on-q}
For $x\in\mathbb{Q}$, define
\[
|x|:=
\begin{cases}
x, & 0_{\mathbb{Q}}\le x,\\
-x, & x<0_{\mathbb{Q}}.
\end{cases}
\]
\end{definition}
\end{definitionbox}

\begin{remark*}[Standard quantified statement]
For every $x\in\mathbb{Q}$,
\[
0_{\mathbb{Q}}\le x\Longrightarrow |x|=x,
\qquad
x<0_{\mathbb{Q}}\Longrightarrow |x|=-x.
\]
\end{remark*}

\begin{remark*}[Interpretation]
The absolute value of a rational number is the number itself when it is
nonnegative and its additive inverse when it is negative.
\end{remark*}

\begin{dependencies}
\begin{itemize}
  \item \hyperref[def:order-on-q]{Order on $\mathbb{Q}$}
  \item \hyperref[def:zero-in-q]{Zero in $\mathbb{Q}$}
  \item \hyperref[def:negation-on-q-classes]{Negation on $\mathbb{Q}$}
\end{itemize}
\end{dependencies}

\begin{theorembox}{Theorem}
\begin{theorem}[Absolute Value Is Nonnegative on $\mathbb{Q}$]
\label{thm:absolute-value-nonnegative-on-q}
For every $x\in\mathbb{Q}$,
\[
0_{\mathbb{Q}}\le |x|.
\]
\hyperref[prf:absolute-value-nonnegative-on-q]{\textit{Go to proof.}}
\end{theorem}
\end{theorembox}

\begin{remark*}[Standard quantified statement]
For every $x\in\mathbb{Q}$,
\[
0_{\mathbb{Q}}\le |x|.
\]
\end{remark*}

\begin{remark*}[Interpretation]
Rational absolute value always produces a nonnegative rational number.
\end{remark*}

\begin{dependencies}
\begin{itemize}
  \item \hyperref[def:absolute-value-on-q]{Absolute Value on $\mathbb{Q}$}
  \item \hyperref[def:order-on-q]{Order on $\mathbb{Q}$}
  \item \hyperref[thm:trichotomy-on-q]{Trichotomy on $\mathbb{Q}$}
\end{itemize}
\end{dependencies}

\begin{theorembox}{Theorem}
\begin{theorem}[Absolute
 Value Vanishes Only at Zero on $\mathbb{Q}$]
\label{thm:absolute-value-zero-iff-on-q}
For every $x\in\mathbb{Q}$,
\[
|x|=0_{\mathbb{Q}}
\quad\Longleftrightarrow\quad
x=0_{\mathbb{Q}}.
\]
\hyperref[prf:absolute-value-zero-iff-on-q]{\textit{Go to proof.}}
\end{theorem}
\end{theorembox}

\begin{remark*}[Standard quantified statement]
For every $x\in\mathbb{Q}$,
\[
|x|=0_{\mathbb{Q}}\Longleftrightarrow x=0_{\mathbb{Q}}.
\]
\end{remark*}

\begin{remark*}[Interpretation]
The only rational number with absolute value zero is zero itself.
\end{remark*}

\begin{dependencies}
\begin{itemize}
  \item \hyperref[def:absolute-value-on-q]{Absolute Value on $\mathbb{Q}$}
  \item \hyperref[thm:absolute-value-nonnegative-on-q]{Absolute Value Is Nonnegative on $\mathbb{Q}$}
  \item \hyperref[thm:trichotomy-on-q]{Trichotomy on $\mathbb{Q}$}
\end{itemize}
\end{dependencies}

\begin{theorembox}{Theorem}
\begin{theorem}[Absolute Value Is Multiplicative on $\mathbb{Q}$]
\label{thm:absolute-value-multiplicative-on-q}
For all $x,y\in\mathbb{Q}$,
\[
|x\cdot y|=|x|\cdot |y|.
\]
\hyperref[prf:absolute-value-multiplicative-on-q]{\textit{Go to proof.}}
\end{theorem}
\end{theorembox}

\begin{remark*}[Standard quantified statement]
For all $x,y\in\mathbb{Q}$,
\[
|x\cdot y|=|x|\cdot |y|.
\]
\end{remark*}

\begin{remark*}[Interpretation]
Rational absolute value converts products to products of absolute values.
\end{remark*}

\begin{dependencies}
\begin{itemize}
  \item \hyperref[def:absolute-value-on-q]{Absolute Value on $\mathbb{Q}$}
  \item \hyperref[thm:positive-multiplication-preserves-order-on-q]{Positive Multiplication
 Preserves Order on $\mathbb{Q}$}
  \item \hyperref[cor:product-of-negatives-on-q]{Product of Negatives on $\mathbb{Q}$}
\end{itemize}
\end{dependencies}

\begin{theorembox}{Theorem}
\begin{theorem}[Triangle Inequality on $\mathbb{Q}$]
\label{thm:triangle-inequality-on-q}
For all $x,y\in\mathbb{Q}$,
\[
|x+y|\le |x|+|y|.
\]
\hyperref[prf:triangle-inequality-on-q]{\textit{Go to proof.}}
\end{theorem}
\end{theorembox}

\begin{remark*}[Standard quantified statement]
For all $x,y\in\mathbb{Q}$,
\[
|x+y|\le |x|+|y|.
\]
\end{remark*}

\begin{remark*}[Interpretation]
The absolute value of a rational sum is bounded above by the sum of the
absolute values.
\end{remark*}

\begin{dependencies}
\begin{itemize}
  \item \hyperref[def:absolute-value-on-q]{Absolute Value on $\mathbb{Q}$}
  \item \hyperref[def:addition-on-q-classes]{Addition on $\mathbb{Q}$}
  \item \hyperref[def:order-on-q]{Order on $\mathbb{Q}$}
  \item \hyperref[thm:q-totally-ordered-set]{$\mathbb{Q}$ Is a Totally Ordered Set}
  \item \hyperref[thm:addition-preserves-order-on-q]{Addition Preserves Order on $\mathbb{Q}$}
\end{itemize}
\end{dependencies}

\begin{theorembox}{Theorem}
\begin{theorem}[Absolute Value of a Reciprocal on $\mathbb{Q}$]
\label{thm:absolute-value-of-reciprocal-q}
For every nonzero $x\in\mathbb{Q}$,
\[
|x^{-1}|=|x|^{-1}.
\]
\hyperref[prf:absolute-value-of-reciprocal-q]{\textit{Go to proof.}}
\end{theorem}
\end{theorembox}

\begin{remark*}[Standard quantified statement]
For every $x\in\mathbb{Q}\setminus\{0_{\mathbb{Q}}\}$,
\[
|x^{-1}|=|x|^{-1}.
\]
\end{remark*}

\begin{remark*}[Interpretation]
Absolute value commutes with taking reciprocals of nonzero rational numbers.
\end{remark*}

\begin{dependencies}
\begin{itemize}
  \item \hyperref[def:absolute-value-on-q]{Absolute Value on $\mathbb{Q}$}
  \item \hyperref[thm:reciprocal-on-q-well-defined]{Reciprocal on $\mathbb{Q}$ Is Well-Defined}
  \item \hyperref[thm:sign-of-reciprocal-q]{Sign of a Reciprocal on $\mathbb{Q}$}
  \item \hyperref[thm:absolute-value-multiplicative-on-q]{Absolute Value Is Multiplicative on $\mathbb{Q}$}
\end{itemize}
\end{dependencies}
